\documentclass{beamer}
\usepackage{amsfonts,amsmath,oldgerm,booktabs}
\usetheme{sintef}
\usepackage{xeCJK}
\IfFileExists{assets/NotoSansCJKsc-Regular.otf}{%
  \setCJKmainfont[Path=assets/,AutoFakeBold]{NotoSansCJKsc-Regular.otf}%
}{%
  \IfFontExistsTF{Noto Serif CJK SC}{%
    \setCJKmainfont{Noto Serif CJK SC}%
  }{%
    \IfFontExistsTF{SimSun}{%
      \setCJKmainfont[AutoFakeBold]{SimSun}%
    }{%
      \setCJKmainfont[BoldFont=FandolSong-Bold.otf]{FandolSong-Regular.otf}%
    }%
  }%
}
\graphicspath{{labreportdata/empirical_progress_20261102_v12/figures/}}

\newcommand{\testcolor}[1]{\colorbox{#1}{\textcolor{#1}{test}}~\texttt{#1}}

\usefonttheme[onlymath]{serif}

\titlebackground*{assets/background}

\newcommand{\hrefcol}[2]{\textcolor{cyan}{\href{#1}{#2}}}

\title{迈向快速且可泛化的开放域决策}
\subtitle{Intern-Decision 作为System-1 决策层}
% \course{Master's Degree in Computer Science}
\author{周家纬、孙哲、赵博涵、张帅宇\\李春晓、孙铮浩、谢嘉庆}
% \IDnumber{1234567}
\date{2026 年 11 月 20 日}

\begin{document}
\maketitle

\begin{frame}{研究问题}
\begin{itemize}
    \item 用 LabQueue 检验 Intern-Decision 0.8B 能否按给定规则选择调度动作，以及在资源规模和输入表示变化时的准确率。
    \item 比较温度校准前后的概率质量，以及置信度区分正确和错误决策的能力。
    \item 在公开 JevBench 上测量思考预算对准确率的影响，并检验按置信度求助的准确率和延迟。
\end{itemize}
\end{frame}

\section{Recall}

\begin{frame}{研究计划概览}
\begin{itemize}
    \item \textbf{实验 1 (E1)：} 复现 Intern Decision 4b，0.8b 的基线表现
    \item \textbf{实验 2 (E2)：} 量“快”值不值，看 Jev 类模型和同参数下LLM的 CoT 回答差距
    \item \textbf{实验 3 (E3)：} 设计置信度机制，看什么时候该求助
    \item \textbf{实验 4 (E4)：} 在jev上微调后，其能力是否可以迁移（以实验室 gpu 任务排班为例），并且验证是否可以容纳多领域内的知识（旧题重测）
\end{itemize}
\end{frame}

\section{Progress}

\begin{frame}[fragile]{E1--E4｜实验环境与自动化脚本}
  \begin{columns}[T]
    \begin{column}{0.35\textwidth}
      \begin{itemize}
        \item \textbf{框架现状:} E1-E4 实验脚本及自动化工作流已构建完成。
        \item \textbf{核心特性:}
        \begin{itemize}
          \item 支持多后端 (HF/XTuner)。
          \item 离线模式 (Offline) 支持。
          \item 开启 GPU 环境后，可直接执行脚本启动任务。
        \end{itemize}
      \end{itemize}
    \end{column}
    \begin{column}{0.65\textwidth}
      \tiny
      \begin{verbatim}
#!/usr/bin/env bash
set -euo pipefail
E1_ROOT="$(cd -- "$(dirname -- "${BASH_SOURCE[0]}")/.." && pwd)"
E1_BACKEND="${1:-hf}"
if [[ "$E1_BACKEND" != hf && "$E1_BACKEND" != xtuner ]]; then
  exit 2
fi
cd "$E1_ROOT"
E1_PYTHON="$E1_ROOT/.venv-$E1_BACKEND/bin/python"
export HF_HUB_OFFLINE=1 TRANSFORMERS_OFFLINE=1 HF_DATASETS_OFFLINE=1
export OMP_NUM_THREADS="${OMP_NUM_THREADS:-8}"
if [[ "$E1_BACKEND" == xtuner ]]; then
  export PYTHONPATH="$E1_ROOT:$E1_ROOT/upstream/Intern-Decision"
fi
"$E1_PYTHON" -m decision_lab.offline_bundle check --backend "$E1_BACKEND"
exec "$E1_PYTHON" scripts/run_e1_baseline.py \
  --backend "$E1_BACKEND" \
  --output "$E1_OUTPUT"
\end{verbatim}
    \end{column}
  \end{columns}
\end{frame}

\begin{frame}{E1 补充｜公开任务的传统监督基线}
  \small
  官方模型的 GPU 复现待运行。传统监督对照使用 TF-IDF + NB、LR、SVM，在三个任务上各训练三个数据种子，共 27 组完整公开验证集评估。

  \medskip
  \centering
  \begin{tabular}{@{}lrr@{}}
    \toprule
    \textbf{任务} & \textbf{模型三种子均值范围} & \textbf{训练多数类基线}\\
    \midrule
    SST-2 & 81.61\%--82.11\% & 50.92\%\\
    RTE   & 49.34\%--51.74\% & 52.71\%\\
    BoolQ & 62.25\%--64.33\% & 62.17\%\\
    \bottomrule
  \end{tabular}

  \medskip
  \raggedright
  \footnotesize SST-2 的三种模型均高于多数类基线，RTE 的三种模型均低于该基线。数据来源为 v12 阶段报告。
\end{frame}

\begin{frame}{E2 补充｜增加思考预算提高 JevBench 准确率}
  \small
  在 JevBench 的 Easy、Original、Hard 共 231 题上，用同一 Qwen3.5-4B 权重比较三档思考预算。触及上限时强制结束思考，再按候选答案后缀的对数概率选答案。

  \medskip
  {\centering
  \setlength{\tabcolsep}{5pt}
  \begin{tabular}{@{}lrrrr@{}}
    \toprule
    \textbf{预算} & \textbf{总体准确率} & \textbf{Hard 准确率} & \textbf{自然结束} & \textbf{原文总体}\\
    \midrule
    关闭思考 & 77.06\% & 58.56\% & --      & 81.4\%\\
    512 token  & 85.71\% & 71.17\% & 89/231  & 84.8\%\\
    2048 token & 91.77\% & 83.78\% & 141/231 & 91.8\%\\
    \bottomrule
  \end{tabular}\par}

  \medskip
  \begin{itemize}
    \item 本次为 \href{https://arxiv.org/abs/2610.02076v3}{LLM-as-Jev 表 16} 的独立部分复现。关闭思考时与原文差约 4.3 个百分点，2048 档的总体准确率接近原文。
    \item 原定 E2 的 632 条请求与 128、512、2048 token 配对实验仍待完成。
  \end{itemize}
\end{frame}

\begin{frame}{E3/E4｜LabQueue-v1 的数据从哪里来？}
  \small
  LabQueue-v1 是我们为实验室 GPU 作业调度自建的合成数据集。脚本用固定随机种子生成模拟场景，四种动作标签均衡。

  \medskip
  \begin{enumerate}
    \item 按目标动作抽样生成集群总容量、当前可用资源、作业需求、分片方案和前置依赖。
    \item 用确定性的 \texttt{solve(state)} 按调度规则计算标准答案，并检查它与抽样时指定的目标动作一致。整个标注过程只执行 Python 规则。
    \item 把同一场景写成两条请求，第二条反转选项顺序。两个版本留在同一数据划分，场景用哈希去重。
  \end{enumerate}
\end{frame}

\begin{frame}{E3/E4｜LabQueue 的四种调度动作}
  \small
  \textbf{模型输入}\quad 调度指令（\texttt{POLICY}）、当前状态（\texttt{state}，含总容量、当前资源、作业需求、分片与依赖信息）和四个顺序可变的候选动作。

  \medskip
  \centering
  \begin{tabular}{@{}ll@{}}
    \toprule
    \textbf{动作} & \textbf{含义}\\
    \midrule
    \texttt{dispatch} & 立即启动完整作业\\
    \texttt{queue} & 保持作业不变，等待调度\\
    \texttt{split} & 按给定方案拆分作业\\
    \texttt{reject} & 拒绝不可行的作业\\
    \bottomrule
  \end{tabular}

  \medskip
  \begin{block}{判分}
    预测动作与标准答案一致得 1 分，否则得 0 分。
  \end{block}
\end{frame}

\begin{frame}{E3/E4｜调度规则的优先级与实际错例}
  \small
  \begin{enumerate}
    \item 有未完成的前置依赖 $\Rightarrow$ \texttt{queue}。
    \item 完整作业超过总容量 $\Rightarrow$ 可以按给定方案分片且每片都能装下时 \texttt{split}，否则 \texttt{reject}。
    \item 完整作业能装下，但当前可用量不足 $\Rightarrow$ \texttt{queue}。
    \item 其余情况 $\Rightarrow$ \texttt{dispatch}。
  \end{enumerate}
  \begin{block}{种子 17 的 IID 测试错例}
    无未完成的前置依赖，当前可用 11 张 GPU 和 98 GB 内存，完整作业只需要 2 张 GPU 和 45 GB 内存。作业可拆分，但资源已经足够，标准答案为 \texttt{dispatch}。模型选了 \texttt{split}，得 0 分。
  \end{block}
\end{frame}

\begin{frame}{E3/E4｜两种分布偏移：资源规模与状态表示}
  \small
  \centering
  \begin{tabular}{@{}lll@{}}
    \toprule
    \textbf{数据划分} & \textbf{资源范围} & \textbf{当前资源的表示}\\
    \midrule
    \texttt{train} / \texttt{test\_iid} & 小 & 直接给可用量\\
    \texttt{test\_ood\_scale} & 大 & 直接给可用量\\
    \texttt{test\_ood\_representation} & 小 & 给占用量、预留量\\
    \texttt{test\_ood} & 大 & 给占用量、预留量\\
    \bottomrule
  \end{tabular}

  \medskip
  \raggedright
  \textbf{可用量}\quad 直接读取 \texttt{available}，或按各资源维度用总量减去占用量和预留量。
  \begin{itemize}
    \item 每个数据种子随机生成 1,200 条训练样本、240 条校准样本和 240 条路由验证样本。
    \item 每种测试条件含 200 个留出场景，每个场景的两个选项顺序变体都在同一划分，共 400 条请求。
  \end{itemize}
\end{frame}

\begin{frame}{E3｜温度校准调整置信度}
  \small
  \begin{columns}[c,onlytextwidth]
    \begin{column}{.49\textwidth}
      \centering
      \includegraphics[width=\linewidth,trim=408bp 46bp 0bp 49bp,clip]{cpu_08b_seed17_test_iid.pdf}\par
      \scriptsize 种子 17 / IID。虚线表示置信度等于实际正确率。
    \end{column}
    \begin{column}{.48\textwidth}
      \textbf{拟合温度}\quad Jev 权重固定。每个种子用 240 条独立校准样本拟合一个 $T$，目标是最小化负对数似然（NLL）。三个种子合计 720 条，拟合三个温度。

      \smallskip
      \textbf{校准置信度}\quad 将输出分数除以 $T$ 再转成概率，让各置信度区间的平均置信度与实际正确率更接近。每个种子的四个留出集共用其温度。

      \smallskip
      \textbf{IID 结果}\quad ECE 从 0.286 降到 0.056，准确率仍为 36.0\%。12 个条件的 ECE 均下降。
    \end{column}
  \end{columns}

  \medskip
  \textbf{同一个错例}\quad 正解是 \texttt{dispatch}，校准前后都选 \texttt{split}，置信度从 54.0\% 降到 31.4\%。
\end{frame}

\begin{frame}{E3｜温度校准前后的四项指标}
  \small
  Intern-Decision 0.8B 在 LabQueue 留出集上的结果。每格为种子 17、43、97 的算术平均，箭头表示原始分数到校准后分数。

  \medskip
  {\centering
  \scriptsize
  \setlength{\tabcolsep}{4pt}
  \begin{tabular}{@{}lcccc@{}}
    \toprule
    \textbf{测试条件} & \textbf{ECE $\downarrow$} & \textbf{NLL $\downarrow$} & \textbf{Brier $\downarrow$} & \textbf{错误 AUROC $\uparrow$}\\
    \midrule
    IID      & $0.268\to0.047$ & $1.751\to1.318$ & $0.820\to0.712$ & $0.599\to0.631$\\
    规模 OOD & $0.262\to0.051$ & $1.773\to1.322$ & $0.820\to0.714$ & $0.591\to0.613$\\
    表示 OOD & $0.338\to0.045$ & $2.192\to1.371$ & $0.914\to0.736$ & $0.550\to0.564$\\
    组合 OOD & $0.362\to0.031$ & $2.275\to1.372$ & $0.943\to0.737$ & $0.550\to0.561$\\
    \bottomrule
  \end{tabular}\par}

  \medskip
  \begin{itemize}
    \item ECE（校准误差）、NLL（负对数似然）和 Brier（概率平方误差）在 12 个种子与测试集组合中均下降。所选动作及准确率保持不变。
    \item 错误 AUROC 衡量低置信度能否排出错题。九个 OOD 条件的 AUROC 差值 95\% 区间均包含零。
  \end{itemize}
\end{frame}

\begin{frame}{E3｜校准后，Jev 更容易挑出错题吗？}
  \small
  \begin{block}{怎样算校准得好？}
    假设 Jev 对 100 道题都报 80\% 的置信度，答对约 80 道就算吻合。若只答对 60 道，就高估了 20 个百分点。ECE 按题数加权平均各置信度区间的误差大小。
  \end{block}

  \begin{itemize}
    \item \textbf{检验目标}\quad 温度校准后，Jev 能否更常把错题排在对题前面？
    \item \textbf{怎么测}\quad 在同一批测试题上，分别用校准前后的置信度从低到高排序。用 AUROC 衡量错题排在对题前面的程度。
    \item \textbf{结果}\quad 九个 OOD 条件的 AUROC 变化，其 95\% 区间均包含零。现有结果尚不能确认校准让 Jev 更容易挑出 OOD 错题。
  \end{itemize}
\end{frame}

\begin{frame}{E3｜LabQueue 的传统分类器路由}
  \small
  \begin{itemize}
    \item \textbf{策略}\quad 结构化逻辑回归先答，低置信度时由梯度提升接手。每个种子的四个测试集共用在验证集定好的门槛。
    \item \textbf{结果}\quad IID 三种子平均准确率从 84.83\% 升至 87.00\%，部分 OOD 条件准确率下降。
    \item LabQueue 已收集 720 条 Intern-Decision 路由验证决策。生成式模型的同题对照仍待完成。
  \end{itemize}
\end{frame}

\begin{frame}{E3｜JevBench 冻结路由纠正 4 道留出题}
  \small
  用 45 题拟合温度 $T=1.585$，另用 45 题定门槛。Intern-Decision 4B 的校准置信度低于门槛时，调用 Qwen3-4B-Thinking-2507。生成无效时保留 S1 答案。

  \medskip
  {\centering
  \setlength{\tabcolsep}{6pt}
  \begin{tabular}{@{}lrrr@{}}
    \toprule
    \textbf{S2 上限} & \textbf{S2 触顶} & \textbf{冻结策略求助} & \textbf{留出准确率}\\
    \midrule
    512 token  & 72/141 & 0/141  & 121/141（85.82\%）\\
    2048 token & 32/141 & 24/141 & 125/141（88.65\%）\\
    \bottomrule
  \end{tabular}\par}

  \medskip
  \begin{itemize}
    \item 2048 档的门槛为 0.7154。24 次求助中纠正 4 题、改错 0 题、回退 16 题。
    \item 相对仅用 S1 提升 2.84 个百分点，按场景组配对的 95\% 区间为 $[0.68,\,5.63]$ 个百分点。同预算随机求助的期望准确率为 86.30\%。
  \end{itemize}
  {\footnotesize 2048 档重跑触顶题，复用 512 档提前结束的回答。两档策略均在同一批 141 道留出题上离线重放。\par}
\end{frame}

\begin{frame}{E3｜HTTP 实测显示准确率与延迟的取舍}
  \small
  141 道留出题，每题执行三种策略，共 423 次本机 HTTP 请求。四个 H200 副本各保持一个在途请求，统计客户端收到并解析结果的时间。

  \medskip
  {\centering
  \setlength{\tabcolsep}{5pt}
  \begin{tabular}{@{}lrrrr@{}}
    \toprule
    \textbf{策略} & \textbf{准确率} & \textbf{P50（秒）} & \textbf{P95（秒）} & \textbf{S2 次数}\\
    \midrule
    仅 S1     & 85.82\% & 0.044 & 0.90   & 0\\
    冻结路由  & 87.94\% & 0.756 & 40.00  & 24\\
    全部 S2   & 78.01\% & 12.11 & 275.88 & 141\\
    \bottomrule
  \end{tabular}\par}

  \medskip
  \begin{itemize}
    \item 冻结路由比仅用 S1 多答对 3 题。24 次求助中有 17 次回退，其中 1 题离线时曾纠正。
    \item S2 升级分支的 P95 为 467.0 秒。延迟口径为预热后的本机 HTTP，跨主机网络和外部 API 尚待测量。
  \end{itemize}
\end{frame}

\begin{frame}{E4｜0.8B 微调前：平均准确率约 36\%--39\%}
  \small
  CPU 推理使用同一套发布权重，LabQueue 微调样本为 0 条。

  \begin{center}
    \begin{tikzpicture}[x=2.65cm,y=.19cm,font=\scriptsize]
      \node[anchor=south west] at (-.58,17.6) {准确率（\%）};
      \foreach \pct in {25,30,35,40} {
        \pgfmathsetmacro{\level}{\pct-24}
        \draw[gray!30] (-.42,\level) -- (3.55,\level);
        \node[anchor=east] at (-.43,\level) {\pct};
      }
      \draw[gray,dashed] (-.42,1) -- (3.55,1);
      \node[above left,gray] at (3.55,1) {均匀随机 25\%};
      \foreach \index/\condition/\first/\second/\third/\average in
        {0/IID/36.00/39.25/38.50/37.92,
         1/规模 OOD/39.25/37.00/40.25/38.83,
         2/表示 OOD/33.75/33.00/40.50/35.75,
         3/组合 OOD/35.75/37.50/36.00/36.42} {
        \fill[blue!70!black] ({\index-.15},{\first-24}) circle (1.65pt);
        \fill[blue!70!black] (\index,{\second-24}) circle (1.65pt);
        \fill[blue!70!black] ({\index+.15},{\third-24}) circle (1.65pt);
        \draw[orange!80!black,very thick] ({\index-.21},{\average-24}) -- ({\index+.21},{\average-24});
        \node[right,orange!80!black] at ({\index+.22},{\average-24}) {\average};
        \node[below] at (\index,0) {\condition};
      }
    \end{tikzpicture}

    \scriptsize 蓝点为种子 17、43、97，橙线为均值。每种子每条件 400 条请求，共 4,800 条。
  \end{center}

  种子 17 的 IID 事后诊断中，仅反转选项顺序就有 86/200 个场景改变了所选动作。
\end{frame}

\begin{frame}{E4｜LabQueue 传统监督基线}
  \small
  \begin{columns}[c,onlytextwidth]
    \begin{column}{.38\textwidth}
      每个种子用 1,200 条规则标签训练六种分类器。

      \medskip
      \textbf{文本输入}\quad 将可见状态写成 JSON，由 TF-IDF 提取词和相邻词对的特征。训练朴素贝叶斯（NB）、逻辑回归（LR）和支持向量机（SVM）。

      \medskip
      \textbf{结构化输入}\quad 直接使用可见数值和布尔字段。训练逻辑回归、随机森林和梯度提升。
    \end{column}
    \begin{column}{.60\textwidth}
      \centering
      \includegraphics[width=.94\linewidth]{baseline_accuracy.pdf}
    \end{column}
  \end{columns}

  \medskip
  \textbf{随机森林与 Jev}\quad IID 为 83.8\% 对 37.9\%，组合 OOD 为 39.3\% 对 36.4\%。

  \smallskip
  随机森林用标签从结构化字段学习选动作。Jev 读文字规则，零样本执行资源比较和优先级判断。微调能缩小多少差距仍待同一测试集上的前后对照。
\end{frame}

\begin{frame}{阶段结论与实验覆盖}
  \small
  \begin{itemize}
    \item \textbf{LabQueue}\quad 0.8B 零样本完成三个种子、十二个留出条件。温度校准改善概率指标，OOD 错题排序的变化仍有不确定性。
    \item \textbf{JevBench}\quad 公开题上的三档思考对照、冻结路由与同机 HTTP 均已完成。更长的生成预算让更多困难题得到有效回答，升级分支仍有明显长尾延迟。
    \item \textbf{尚待完成}\quad E1 的官方七套件与 pilot，E2 的完整固定样本对照，E3 的 LabQueue 生成式路由及跨节点 API，E4 的三种子微调与能力保持评测。
  \end{itemize}
\end{frame}

\section{Supplementary}

\newcommand{\cpucondition}[2]{%
  \begin{frame}{0.8B CPU 补充实测：#1}
    \centering
    \includegraphics[width=\textwidth,height=.73\textheight,keepaspectratio]{cpu_08b_seed17_#2.pdf}\par
    \smallskip
    \footnotesize 种子 17。左图为同一批 400 条请求的动作准确率，右图为 0.8B 校准前后的可靠性。
  \end{frame}
}

\cpucondition{IID}{test_iid}
\cpucondition{规模 OOD}{test_ood_scale}
\cpucondition{表示 OOD}{test_ood_representation}
\cpucondition{组合 OOD}{test_ood}

\begin{frame}{0.8B 错误检测：按场景组估计 AUROC 的不确定性}
  \small
  \begin{columns}[c,onlytextwidth]
    \begin{column}{.70\textwidth}
      \centering
      \includegraphics[width=\linewidth,height=.79\textheight,keepaspectratio,trim=0bp 39bp 0bp 51bp,clip]{cpu_error_auroc_intervals.pdf}
    \end{column}
    \begin{column}{.28\textwidth}
      每个条件对 200 个场景组做 1,000 次配对重采样，同一场景的两个变体一起抽取。

      \medskip
      右图给出校准前后差值。九个 OOD 区间均包含零。

      \medskip
      \scriptsize 事后分析，区间未经多重比较校正。沿用原有 S1 预测，尚无 S2 结果。
    \end{column}
  \end{columns}
\end{frame}

\begin{frame}{分类器冻结路由：修正、损害和升级比例}
  \centering
  \includegraphics[width=\textwidth,height=.71\textheight,keepaspectratio]{fix_harm_and_budget.pdf}\par
  \smallskip
  \footnotesize 结构化逻辑回归 $\to$ 梯度提升。25\% 是选阈值时的验证集预算，测试集升级比例随数据分布变化。
\end{frame}

\begin{frame}{分类器路由：同样的升级次数，怎么挑题？}
  \centering
  \includegraphics[width=\textwidth,height=.71\textheight,keepaspectratio]{matched_budget_routing.pdf}\par
  \smallskip
  \footnotesize 横轴为升级比例，纵轴为路由后的准确率。每个预算内各策略升级数量相同，虚线是使用测试标签的理论上界。
\end{frame}

\backmatter
\end{document}
